\section{Anilisis Transitorio}
En esta parte del trabajo vamos a analizar el esquemático que se puede ver en la Figura \ref{fig:circuito_tran}. 
\begin{figure}[H]
    \centering
    \begin{tikzpicture}[transform shape]
	% Paths, nodes and wires:
	\node[spdt] at (8.595, 1){};
	\draw (9.19, 0.685) -| (12, 0.5) -- (12, 0.25);
	\draw (12, 0.25) to[american inductor, l_={$L$}] (12, -1.75);
	\draw (12, -1.75) to[american resistor, l_={$R_3$}] (12, -3.5);
	\draw (10, 0.75) to[capacitor, l_={$C$}] (10, -1.75);
	\draw (10, -1.75) to[american resistor, l_={$R_2$}] (10, -3.5);
	\draw (5.75, 1) to[american resistor, l_={$R_f$}] (5.75, -1.5);
	\draw (5.75, -1.5) to[battery1, l_={$V_i$}] (5.75, -3.5);
	\draw (5.75, -3.5) -- (12, -3.5);
	\node[rground] at (6.25, -3.5){};
	\draw (5.75, 1) to[american resistor, l={$R_1$}] (8, 1);
	\draw[line width=0.6pt, -latex] (12.75, -3.5) -- (12.75, 0.5);
	\node[shape=rectangle, minimum width=0.715cm, minimum height=0.465cm](N1) at (13.375, -1.5){} node[anchor=center] at (N1.text){$V_0$};
\end{tikzpicture}
    \caption{Circuito utilizado para el análisis del transitorio}
    \label{fig:circuito_tran}
\end{figure}

A su vez los datos utilizados para el trabajo se encuentran en la Tabla \ref{tab:datos_transitorio}. 

\begin{table}[H]
    \centering
    \begin{tabular}{|c|c|c|}
        \hline
        Elemento & Valor nominal & Valor medido \\
        \hline
        $V_i$ & $10\,\text{V}$ & \\
        \hline
        $R_1$ & $100\,\Omega$ & $98\,\Omega$ \\
        \hline
        $R_2\, [\Omega]$ & $8{,}2$ & $8{,}4$ \\
        \hline
        $R_3 \,[\Omega]$ & $6{,}8$ & $7{,}2$ \\
        \hline
        $R_f\,[\Omega]$ & $50$ & \\
        \hline
        $L\,[\text{mH}]$ & $1$ & $0{,}988$ \\
        \hline
        $C\,[\text{nF}]$ & $100$ & $107$ \\
        \hline
    \end{tabular}
    \caption{Valores nominales y medidos de los elementos utilizados en el análisis transitorio}
    \label{tab:datos_transitorio}
\end{table}

\subsection{Analisis de la carga}
Para esta etapa 

\subsection{Analisis de la descarga}

\subsection{Circuito sin resistencias}

Si se analiza idealmente el circuito de la Figura \ref{fig:circuito_tran} suponiendo que $R_1=R_2=R_3=R_f=0\,\Omega$, el circuito no tendría ningún mecanismo de disipación de energía. En la etapa de carga, al cerrar la llave, la fuente ideal quedaría aplicada directamente sobre el capacitor y la bobina. Como el capacitor inicialmente descargado no puede cambiar su tensión instantáneamente sin una corriente impulsiva, el modelo ideal predice un pico de corriente teóricamente infinito en el instante inicial.

Además, al quedar la bobina ideal conectada directamente a una fuente continua ideal, su corriente crecería sin límite con el tiempo, ya que se cumple:

\[
V_L = L \frac{di_L}{dt}
\]

Por lo tanto, si $V_L$ es constante y distinto de cero, la corriente por la bobina aumenta linealmente.

En cambio, durante la descarga, si la fuente queda desconectada y solo permanece el conjunto ideal $LC$, la energía almacenada se intercambiaría indefinidamente entre el campo eléctrico del capacitor y el campo magnético de la bobina. En ese caso aparecería una oscilación senoidal no amortiguada, con frecuencia natural:

\[
\omega_0 = \frac{1}{\sqrt{LC}}
\]

Este comportamiento no puede ocurrir exactamente en la realidad, ya que todos los elementos poseen resistencias parásitas y pérdidas asociadas. Por eso, en un circuito real, la oscilación se amortigua y la energía almacenada termina disipándose en forma de calor.

\subsection{Diagrama de montecarlo}